% Chapter 1: The Briefing
% Revised 2026-04-09: integrated humor passages, applied trim notes.
% Scene function: Tert receives the case file. Voice, world, and the
% motivation gap establish simultaneously. Gordon introduced bureaucratically.
% Turn: Tert's recognition of the living room's acoustic canvas.

\chapter{On Paper}

On paper, it is not a remarkable assignment.

The file is on my desk when I come in. Someone has dropped it there during the off-shift, which means Intake has a preference about when I look at it and does not want to say so in writing. I note the preference and put the file where I would have put it anyway, in the middle of the blotter, and sit down.

The lights in this corridor of the office hum at a frequency that is not quite fluorescent and not quite anything else. Facilities has been notified. The lights have not changed.

The file is standard Intake format. Cover sheet, subject profile, environmental survey, proposed escalation schedule, performance benchmarks, my signature block in the lower right corner where I am expected to agree that I have read the file and accept the case. I do not sign anything yet. Signing is the last thing you do. If you sign at the front of the process, Management learns to expect it, and then they stop respecting the distinction between read and accepted, and you wind up signing for things you have not read. I learned this once. I did not need to learn it twice.

I flip to the subject profile.

Subject: male, 55, widower. Retired. Sole resident of a single-family dwelling. Religious affiliation: none on record. Church attendance: none on record. Recent significant loss: spouse, deceased 10 months, ALS.

I look at the religious affiliation line for a while.

The stated objective of any atmospheric haunting, per the Operations Manual (seventh edition, Chapter 4, Section 2, subsection ``Purpose''), is to erode the subject's spiritual defenses through sustained environmental dread, leading to a crisis of faith, leading to rejection of divine protection, leading to a harvestable soul. That is the pipeline. That is what we are for.

My subject does not appear to have spiritual defenses. He does not attend church. He has no religious affiliation. He is not, as far as Intake can tell, protected by anything except a deadbolt and a recently serviced smoke detector.

There is no field on the intake form for ``subject does not believe in the thing we need him to reject.'' There is a field for ``low religious engagement,'' which is what Intake has checked, and which is supposed to indicate an easier case. Low defenses. Quick timeline. High likelihood of favorable Soul Outcome. I am not sure how you produce a favorable Soul Outcome with a subject who does not appear to have a relationship with God to sever, but the form does not have a field for that either, and this is not a problem I am going to solve at my desk on a Tuesday morning.

In the notes column someone has written, in Intake shorthand, \emph{caregiver background; expect low avoidance, high baseline fatigue.}

I reread the caregiver note. I appreciate whoever wrote it. Most Intake officers do not think about baseline fatigue. They look at the loss column and assume grief does the work for them. They have never worked a caregiver subject. A caregiver subject has already spent three years learning to sleep through terror. They have been woken up at three in the morning by real things. A floor creak is not terror to them. A floor creak is a normal Tuesday.

I make a note in my margin pad: \emph{adjust Phase 1 expectations. Ambient floor will need to be below normal-Tuesday level before any wrongness reads at all.}

I flip to the proposed escalation schedule. Intake has pre-filled it from the standard template: Week 1, environmental conditioning. Week 2, localized manifestation. Week 3, direct interface. Week 4, escalated confrontation. Week 5, resolution pressure.

I have seen this template four hundred times. I do not know who wrote it. I suspect it was written by someone who has never conducted a field haunting, because ``direct interface'' in week three is aggressive by any reasonable standard. Direct interface in week three is what Apparitions does. Apparitions shows up, makes a face, and leaves. It is the equivalent of shouting at someone in a parking lot and calling it persuasion.

I cross out the week three entry and write in \emph{sustained ambient pressure; consolidate Phase 1 and 2 gains.} I do this with every case file. Nobody has ever noticed. I am beginning to suspect nobody reads these after I sign them.

I flip to performance benchmarks. Fear Index target: 7.2 (sustained). Timeline adherence: 90\%. Collateral Awareness: keep below 15\%. Soul Outcome: positive.

I do not know what a Fear Index of 7.2 means. I do not know what a Fear Index of 7.2 is measured against. I have asked. The answer involves a calibration document that references a second calibration document that references a regional standards meeting from 1987 that nobody has the minutes for. Somewhere in that chain of references is presumably a human being who was exactly 7.2 afraid, against whom all subsequent fear is measured, and I would like to meet them, because I have questions.

I sign it.

I flip to environmental.

The dwelling is a single-story bungalow, 1923, original hardwood, original plaster, original windows (storm-upgraded), an attic accessed by a pull-down ladder in the hallway, a basement with an exterior entrance, and a back yard with a garden that has not been tended in ten months but, per the surveyor's note, still has structure. Three bedrooms, one converted to a home office, one converted to a sickroom and not yet reconverted. Acoustic survey attached.

I do not open the acoustic survey yet. I am being disciplined. I am reading the file in order.

The sickroom line catches me because sickrooms are environmental features I know how to work with. They do not require supernatural intervention. They require letting the room do what the room already wants to do. The adjustment furniture generates its own wrongness. A grab bar on a hallway wall does not belong there and the human eye knows it even when the human mind has forgotten. A raised toilet seat in a bathroom without a patient is a loud object. I have worked sickroom subjects before. They go quickly.

But this one is ten months post. Most of the medical equipment will be gone by now. The subject will have packed most of it up or donated it. What remains will be the things the subject could not bring himself to move. Those are the ones I want.

The window behind my desk shows something different today. I do not look at it. I am looking at photographs.

I flip to the environmental survey's photograph insert.

There are four photographs. Two exterior, two interior. I look at the exteriors first. The house is painted a color I do not have a name for, between gray and green, with cream trim. Front porch with a swing. A hedge along one side. A maple in the front yard that is bare in the photograph, which means the survey was done in winter. The house sits back from a quiet street in a grid of similar houses. Old neighborhood. The kind of street that hasn't had a new house built on it in eighty years.

I look at the interior photographs.

The first is a living room. Hardwood floor, high ceiling for a bungalow, a fireplace with a mantel, a pair of armchairs facing slightly away from each other in a configuration that means they were once facing each other and someone has moved one, built-in bookshelves on either side of the fireplace, a standing lamp, a braided rug, a coffee table with a coaster but no cup. On the mantel, a photograph I cannot see the detail of. Above the mantel, an empty nail.

The second is a hallway. Hardwood, runner rug, a grab bar installed on the left-hand wall at about hip height, bolted into a stud. A doorway open at the end of the hall showing a bed made with a white quilt. A framed print on the wall too small to make out. A light fixture on the ceiling I can make out well enough to know that it is original to the house and has not been rewired since roughly the Truman administration.

I sit with the photographs for a moment.

Then I go back to the first one and look at the bookshelves.

Built in, flanking the fireplace, floor to ceiling. Even in the survey photograph, even at this resolution, I can count the shelves. Six on each side. And I can tell that they are full and that the books are arranged by a person, not by an aesthetic. I can tell because the bindings do not match. Paperbacks and hardcovers together. A few reference volumes on the bottom shelves, the way someone puts the heavy books down where the weight belongs. On one of the shelves, a book is lying flat on top of the other books, which is what happens when a book is being read and there is no bookmark handy and someone puts it down in a hurry meaning to come back.

I do not know what the book is. That is not what I am looking at.

I am looking at the acoustics.

A bookshelf that tall, full of that many books, in a room with a hardwood floor and plaster walls and a fireplace, is a room that has been accidentally tuned. Books dampen high frequencies and let low frequencies through. A room with a wall of books on each side of a central feature does not echo. It absorbs the top end. Which means anything I do in that room at the bass end, anything below about 200 hertz, anything at the register of a house settling, will be audible without being locatable. It will arrive everywhere at once. The subject will not be able to point to a direction and say \emph{the sound is coming from over there.} The sound will be coming from the room.

I have been waiting my entire career for a room like this.

I sign the acceptance block at the bottom of the cover sheet. My designation goes on the signature line. Tertius Gradus, Atmospheric Hauntings, Operative 47-C. Colleagues call me Tert. Management calls me by whichever form of address gets them the response they want. I am not in the habit of correcting either.

I walk the signed file over to Intake's pickup window. There is no one at the window. There never is. I put the file in the tray and I go back to my desk and I open my own calendar and I block out the first five weeks of a standard atmospheric cycle.

Then I sit for a moment looking at the calendar, at the five weeks stretched across it in clean bureaucratic gray, and I think about bookshelves.
